\documentclass[11pt,a4paper]{report}

% Packages de base
\usepackage[utf8]{inputenc}
\usepackage[T1]{fontenc}
\usepackage{lmodern}
\usepackage[french]{babel}
\usepackage{geometry}
\geometry{a4paper, margin=2.5cm}
\usepackage{graphicx}
\usepackage{amsmath, amssymb}
\usepackage{array}
\usepackage{booktabs} % Pour de jolis tableaux
\usepackage{enumitem} % Pour personnaliser les listes
\usepackage{hyperref} % Pour les liens cliquables
\hypersetup{
    colorlinks=true,
    linkcolor=blue,
    filecolor=magenta,
    urlcolor=cyan,
    pdftitle={Rapport de Stage - Housy},
    pdfpagemode=FullScreen,
}

% Packages pour les couleurs et les boîtes colorées
\usepackage[svgnames,table]{xcolor} % xcolor avec options svgnames et table pour les tableaux
\usepackage{tcolorbox}
\tcbuselibrary{skins,breakable,hooks} % Pour des tcolorbox plus avancées

% Définition de couleurs personnalisées (palette douce et professionnelle)
\definecolor{HousyBlue}{HTML}{2A628F} % Bleu principal
\definecolor{HousyGreen}{HTML}{4CAF50} % Vert pour succès, points positifs
\definecolor{HousyOrange}{HTML}{FF9800} % Orange pour avertissements, tips
\definecolor{HousyLightBlue}{HTML}{E3F2FD} % Bleu clair pour fonds de boîtes
\definecolor{HousyLightGreen}{HTML}{E8F5E9} % Vert clair
\definecolor{HousyLightOrange}{HTML}{FFF3E0} % Orange clair
\definecolor{HousyGrey}{HTML}{757575} % Gris pour texte secondaire
\definecolor{HousyDarkBlue}{HTML}{1A237E} % Bleu foncé pour titres importants

% Configuration des titres de section avec couleurs
\usepackage{titlesec}
\titleformat{\chapter}[display]
  {\normalfont\huge\bfseries\color{HousyDarkBlue}}
  {\chaptertitlename\ \thechapter}{20pt}{\Huge}
\titleformat{\section}
  {\normalfont\Large\bfseries\color{HousyBlue}}
  {\thesection}{1em}{}
\titleformat{\subsection}
  {\normalfont\large\bfseries\color{HousyGreen}}
  {\thesubsection}{1em}{}
\titleformat{\subsubsection}
  {\normalfont\normalsize\bfseries\color{HousyOrange}}
  {\thesubsubsection}{1em}{}

% Style pour les légendes de figures et tableaux
\usepackage[labelfont={bf,color=HousyBlue},textfont=it]{caption}

% Environnement pour les "tips" ou points clés
\newtcolorbox{tipbox}[1][]{
  breakable,
  colback=HousyLightOrange!80!white, % Couleur de fond légèrement orangée
  colframe=HousyOrange, % Couleur du cadre
  fonttitle=\bfseries,
  coltitle=black, % Couleur du titre de la boîte
  title=#1,
  arc=4mm, % Coins arrondis
  boxrule=1pt, % Épaisseur du cadre
  enhanced,
  shadow={2mm}{-2mm}{0mm}{black!20!white} % Ombre portée
}

% Environnement pour les encadrés d'information
\newtcolorbox{infobox}[1][]{
  breakable,
  colback=HousyLightBlue!70!white, % Couleur de fond bleu clair
  colframe=HousyBlue, % Couleur du cadre
  fonttitle=\bfseries,
  coltitle=black,
  title=#1,
  arc=4mm,
  boxrule=1pt,
  enhanced,
  shadow={2mm}{-2mm}{0mm}{black!20!white}
}

% Environnement pour le code (simple pour l'exemple, listings serait mieux pour du vrai code)
\usepackage{verbatim}
\newenvironment{codeblock}[1][Code Example]
  {\begin{tcolorbox}[colback=black!5!white, colframe=HousyGrey!70!black, title=#1, fonttitle=\bfseries\small, arc=2mm, boxrule=0.5pt,breakable]}
  {\end{tcolorbox}}

% Configuration pour les tableaux colorés
\usepackage{colortbl} % Pour colorer les cellules, lignes, colonnes
\arrayrulecolor{HousyBlue} % Couleur des lignes du tableau
\setlength{\arrayrulewidth}{0.5pt} % Épaisseur des lignes

\begin{document}

% Supposons que ce soit le chapitre 3
\chapter{Conception et Réalisation Technique de Housy}
\label{chap:conception_realisation}

Ce chapitre central détaille l'architecture technique, les choix technologiques, ainsi que l'organisation et les fonctionnalités clés des différentes composantes du projet Housy. Il vise à fournir une compréhension approfondie du système développé, de sa base de données à son interface utilisateur, en passant par la logique métier implémentée côté serveur.

\section{Architecture et Analyse Technique de Housy}
\label{sec:architecture_analyse}
L'architecture de Housy a été pensée pour être modulaire, évolutive et maintenable, en s'appuyant sur des technologies modernes et éprouvées. Cette section décrit la structure globale du système et les choix technologiques qui la sous-tendent.

\subsection{Vue d’ensemble du Système}
\label{ssec:vue_ensemble}
Housy est une application web fullstack conçue pour [Objectif principal de Housy, par exemple: la gestion immobilière intelligente / la mise en relation d'acheteurs et vendeurs / etc.]. Elle s'articule autour d'une architecture client-serveur classique, avec une séparation claire des préoccupations entre le frontend, le backend et la base de données.

\begin{figure}[h!]
    \centering
    \includegraphics[width=0.8\linewidth]{./images/placeholder_architecture_globale.png} % Remplacez par le chemin réel
    \caption{Diagramme de l'architecture globale de Housy.}
    \label{fig:architecture_globale}
\end{figure}
\textit{Le diagramme ci-dessus illustre les principaux composants de Housy et leurs interactions.}

L'utilisateur interagit avec l'application via une interface web (frontend) développée en React.js. Cette interface communique avec un serveur backend (développé en Node.js avec Express.js) via une API RESTful. Le backend gère la logique métier, l'authentification, les autorisations et les interactions avec la base de données PostgreSQL, qui stocke toutes les données persistantes de l'application. Des services tiers peuvent également être intégrés, par exemple pour la géolocalisation, les paiements, ou l'envoi d'emails.

\begin{infobox}{Communication Client-Serveur}
La communication entre le frontend et le backend est principalement effectuée via des requêtes HTTP (GET, POST, PUT, DELETE) suivant les principes REST. Les données sont échangées au format JSON. Pour des fonctionnalités temps-réel (par exemple, notifications), des WebSockets pourraient être envisagés dans une future itération.
\end{infobox}

\subsection{Stack Technique et Choix des Technologies}
\label{ssec:stack_technique}
Le choix des technologies pour Housy a été guidé par plusieurs facteurs : la performance, la scalabilité, la disponibilité des compétences, la richesse de l'écosystème (bibliothèques, communauté) et la facilité de développement.

\begin{table}[h!]
\centering
\caption{Stack technique principale de Housy.}
\label{tab:stack_technique}
\rowcolors{2}{HousyLightBlue!50}{white} % Alternance de couleurs pour les lignes
\begin{tabular}{>{\bfseries}l p{0.6\linewidth}}
\toprule
\rowcolor{HousyBlue!80} % Couleur pour l'en-tête
\textcolor{white}{Composant} & \textcolor{white}{Technologies / Outils} \\
\midrule
Frontend & React.js, Redux (ou Zustand/Context API), Next.js (pour SSR/SSG si pertinent), HTML5, CSS3 (avec SASS/Styled Components), Axios \\
Backend & Node.js, Express.js, TypeScript, JWT (JSON Web Tokens) pour l'authentification, Sequelize (ORM) \\
Base de Données & PostgreSQL \\
Tests & Jest, React Testing Library (Frontend), Mocha/Chai/Supertest (Backend) \\
Déploiement & Docker, Nginx, CI/CD (GitLab CI, GitHub Actions, Jenkins) \\
Outillage & Git, VS Code, ESLint, Prettier, Postman (ou Insomnia) \\
IA (si applicable) & [Ex: Python, TensorFlow/PyTorch, Scikit-learn, API de services IA Cloud] \\
\bottomrule
\end{tabular}
\end{table}
\textit{Le tableau ci-dessus synthétise les technologies clés utilisées dans le projet Housy.}

\begin{tipbox}{Justification des choix technologiques}
\textbf{React.js (Frontend)}: Choisi pour sa popularité, son écosystème riche, son approche basée sur les composants facilitant la réutilisabilité et la maintenance, et sa performance grâce au DOM virtuel.
\textbf{Node.js/Express.js (Backend)}: Permet d'utiliser JavaScript/TypeScript sur l'ensemble de la stack (cohérence), offre d'excellentes performances pour les applications I/O intensives et dispose d'une vaste communauté et de nombreux modules (NPM).
\textbf{PostgreSQL (Base de Données)}: Base de données relationnelle open-source robuste, fiable, et très respectueuse des standards SQL. Elle offre des fonctionnalités avancées comme le support JSONB, les transactions ACID, et une bonne scalabilité.
\textbf{TypeScript}: Utilisé à la fois pour le frontend et le backend afin d'améliorer la robustesse du code grâce au typage statique, de faciliter la refactorisation et la collaboration au sein de l'équipe.
\end{tipbox}

\subsection{Frontend : Organisation et Fonctionnalités}
\label{ssec:frontend}
Le frontend de Housy est responsable de l'interface utilisateur et de l'expérience utilisateur. Il est développé en React, en suivant une architecture basée sur les composants.

\subsubsection{Structure des Composants et Organisation du Code}
Le code source du frontend est structuré de manière modulaire pour faciliter la navigation, la compréhension et la maintenance. Une structure typique pourrait être :
\begin{codeblock}{Exemple de structure de dossiers Frontend (React)}
src/
|-- assets/          # Images, polices, etc.
|-- components/      # Composants UI réutilisables (atomes, molécules)
|   |-- common/      # Boutons, Inputs, Modales...
|   |-- layout/      # Header, Footer, Sidebar...
|-- features/        # Modules fonctionnels (ex: Annonces, ProfilUtilisateur)
|   |-- annonces/
|   |   |-- components/ # Composants spécifiques à la feature Annonces
|   |   |-- hooks/      # Hooks spécifiques
|   |   |-- services/   # Appels API pour les annonces
|   |   |-- AnnoncesList.tsx
|   |   |-- AnnonceDetail.tsx
|-- hooks/           # Hooks custom globaux (useAuth, useTheme...)
|-- pages/           # Composants de haut niveau mappés aux routes (vues)
|-- services/        # Services globaux (api.ts, authService.ts)
|-- store/           # Gestion d'état global (Redux, Zustand, Context)
|   |-- slices/      # Reducers/Slices pour Redux Toolkit
|   |-- index.ts
|-- styles/          # Styles globaux, thèmes
|-- utils/           # Fonctions utilitaires
|-- App.tsx          # Composant racine
|-- index.tsx        # Point d'entrée
\end{codeblock}
\textit{L'organisation des fichiers et dossiers du frontend vise la modularité et la séparation des préoccupations.}

\begin{figure}[h!]
    \centering
    \includegraphics[width=0.7\linewidth]{./images/placeholder_composants_ui.png}
    \caption{Exemple de décomposition d'une page Housy en composants React.}
    \label{fig:composants_ui}
\end{figure}
\textit{Une capture d'écran ou un schéma illustrant la hiérarchie des composants pour une page clé de Housy serait pertinent ici.}

\subsubsection{Routing, Gestion d’État, et Flux de Données}
Le \textbf{routing} est géré par `react-router-dom`, permettant de naviguer entre les différentes vues de l'application sans recharger la page.
La \textbf{gestion d'état} est un aspect crucial. Pour les états locaux des composants, les hooks `useState` et `useReducer` de React sont utilisés. Pour l'état global de l'application (par exemple, informations de l'utilisateur connecté, préférences), une solution comme Redux Toolkit (avec ses `slices` et son `store` centralisé) ou Zustand (plus léger) est implémentée. L'API Context de React peut être utilisée pour des états partagés moins complexes.

\begin{infobox}{Gestion d'état avec Redux Toolkit}
Redux Toolkit simplifie l'écriture du code Redux en fournissant des abstractions comme `createSlice` pour générer automatiquement les actions et les reducers, et `configureStore` pour une configuration optimisée du store. Cela réduit le boilerplate et améliore la maintenabilité.
\end{infobox}

\subsubsection{Accessibilité (a11y) et Responsive Design}
L'\textbf{accessibilité} est prise en compte dès la conception en utilisant des balises HTML sémantiques, des attributs ARIA lorsque nécessaire, et en s'assurant du contraste des couleurs et de la navigabilité au clavier.
Le \textbf{responsive design} est assuré via des Media Queries CSS, des grilles flexibles (Flexbox, CSS Grid) et potentiellement des bibliothèques UI comme Material-UI ou Bootstrap (si utilisées) qui offrent des composants responsive par défaut. L'objectif est que Housy soit utilisable sur une large gamme d'appareils (ordinateurs, tablettes, smartphones).

\subsubsection{Points d’Optimisation et Sécurité Côté Client}
\textbf{Optimisations Frontend :}
\begin{itemize}[label=\textcolor{HousyGreen}{\textbullet}]
    \item \textbf{Code Splitting} : Avec React.lazy et Suspense (ou via Next.js), pour ne charger que le code JavaScript nécessaire à la vue actuelle.
    \item \textbf{Memoization} : Utilisation de `React.memo` pour les composants, `useMemo` et `useCallback` pour éviter les re-renderings inutiles.
    \item \textbf{Optimisation des images} : Compression, formats modernes (WebP), lazy loading.
    \item \textbf{Minification et compression} : Des assets (JS, CSS) lors du build.
    \item \textbf{Caching} : Utilisation des stratégies de cache du navigateur et potentiellement des Service Workers pour une PWA.
\end{itemize}

\textbf{Sécurité Côté Client :}
\begin{itemize}[label=\textcolor{HousyOrange}{\textbullet}]
    \item \textbf{Validation des entrées} : Toujours valider les données côté client pour une meilleure UX, mais \textbf{impérativement re-valider côté serveur}.
    \item \textbf{Protection contre XSS (Cross-Site Scripting)} : React échappe par défaut les données affichées dans le JSX. Éviter `dangerouslySetInnerHTML` sans sanitization.
    \item \textbf{Gestion des Tokens JWT} : Stockage sécurisé des tokens (par exemple, dans des cookies HttpOnly si possible, ou localStorage avec des mesures additionnelles).
    \item \textbf{HTTPS} : L'application doit être servie exclusivement via HTTPS.
\end{itemize}
\textit{Une capture d'écran d'un outil d'audit de performance (Lighthouse) ou d'un exemple de code d'optimisation pourrait être insérée ici.}

\subsection{Backend : Organisation et Fonctionnalités}
\label{ssec:backend}
Le backend de Housy, développé en Node.js avec le framework Express.js et TypeScript, est le cœur de l'application. Il gère la logique métier, l'authentification, les interactions avec la base de données et expose une API RESTful au frontend.

\subsubsection{Architecture Serveur et API}
L'architecture du backend suit souvent un modèle en couches (Layered Architecture) ou un modèle N-tiers pour une meilleure séparation des préoccupations :
\begin{itemize}
    \item \textbf{Couche de Présentation/Route} : Gère les requêtes HTTP entrantes, extrait les paramètres, et appelle les services appropriés. Définit les endpoints de l'API. (Ex: `routes/annonces.routes.ts`)
    \item \textbf{Couche de Service/Logique Métier} : Contient la logique applicative principale. Orchestre les opérations, valide les données métier, et interagit avec la couche d'accès aux données. (Ex: `services/annonce.service.ts`)
    \item \textbf{Couche d'Accès aux Données (DAL)} : Responsable de la communication avec la base de données, souvent via un ORM comme Sequelize. (Ex: `models/annonce.model.ts` et les opérations CRUD associées).
\end{itemize}

\begin{codeblock}{Exemple de structure de dossiers Backend (Node.js/Express/TypeScript)}
src/
|-- config/          # Fichiers de configuration (DB, JWT secret, ports...)
|-- controllers/     # (Ou Handlers) Logique de gestion des requêtes/réponses HTTP
|   |-- annonce.controller.ts
|   |-- auth.controller.ts
|-- DTOs/            # Data Transfer Objects, validation des données d'entrée/sortie
|   |-- create-annonce.dto.ts
|-- middlewares/     # Middlewares Express (auth, error handling, logging...)
|   |-- auth.middleware.ts
|   |-- error.middleware.ts
|-- models/          # Définitions des modèles de données (ex: Sequelize models)
|   |-- annonce.model.ts
|   |-- user.model.ts
|-- routes/          # Définition des routes de l'API
|   |-- index.routes.ts  # Routeur principal
|   |-- annonce.routes.ts
|   |-- auth.routes.ts
|-- services/        # Logique métier, interaction avec les modèles
|   |-- annonce.service.ts
|   |-- user.service.ts
|-- utils/           # Fonctions utilitaires (logger, helpers...)
|-- app.ts           # Configuration de l'application Express
|-- server.ts        # Point d'entrée du serveur HTTP
\end{codeblock}
\textit{L'organisation des fichiers et dossiers du backend est cruciale pour la maintenabilité et l'évolutivité.}

\begin{figure}[h!]
    \centering
    \includegraphics[width=0.8\linewidth]{./images/placeholder_backend_layers.png}
    \caption{Schéma illustrant l'architecture en couches du backend Housy.}
    \label{fig:backend_layers}
\end{figure}
\textit{Un schéma de l'architecture en couches du backend, montrant le flux d'une requête, serait approprié ici.}

\subsubsection{Endpoints, Middlewares, Sécurité et Gestion des Erreurs}
\textbf{Endpoints API :} L'API est conçue selon les principes RESTful. Par exemple :
\begin{itemize}
    \item `GET /api/annonces` : Lister toutes les annonces.
    \item `POST /api/annonces` : Créer une nouvelle annonce.
    \item `GET /api/annonces/:id` : Obtenir une annonce spécifique.
    \item `PUT /api/annonces/:id` : Mettre à jour une annonce.
    \item `DELETE /api/annonces/:id` : Supprimer une annonce.
    \item `POST /api/auth/login` : Connexion utilisateur.
    \item `POST /api/auth/register` : Inscription utilisateur.
\end{itemize}

\textbf{Middlewares :} Express.js utilise largement les middlewares pour traiter les requêtes. Exemples courants dans Housy :
\begin{itemize}
    \item `express.json()` : Pour parser les corps de requête JSON.
    \item `cors` : Pour gérer les Cross-Origin Resource Sharing.
    \item Middleware d'authentification : Vérifie la validité des tokens JWT pour protéger les routes.
    \item Middleware de logging : Enregistre les informations sur les requêtes.
    \item Middleware de gestion des erreurs : Centralise la capture et le formatage des erreurs.
\end{itemize}

\begin{tipbox}{Sécurité Backend}
La sécurité est primordiale :
\begin{itemize}
    \item \textbf{Authentification robuste} : Utilisation de JWT, mots de passe hashés (bcrypt).
    \item \textbf{Autorisation} : Vérifier que l'utilisateur a les droits nécessaires pour accéder à une ressource ou effectuer une action (RBAC - Role-Based Access Control).
    \item \textbf{Validation des entrées} : Toujours valider et assainir toutes les données provenant du client (côté serveur) pour prévenir les injections (SQL, NoSQL), XSS (si le backend génère du HTML), etc. Utilisation de bibliothèques comme `class-validator` avec les DTOs.
    \item \textbf{Protection contre les attaques courantes} : CSRF (si utilisation de cookies pour sessions), DoS/DDoS (limitation de débit), etc. Utilisation de `helmet` pour sécuriser les en-têtes HTTP.
    \item \textbf{Gestion des secrets} : Ne jamais stocker de secrets (clés API, mots de passe DB) dans le code. Utiliser des variables d'environnement et des outils de gestion de secrets.
    \item \textbf{Logging et monitoring} : Pour détecter et réagir aux activités suspectes.
\end{itemize}
\end{tipbox}

\textbf{Gestion des Erreurs :} Un middleware centralisé de gestion des erreurs est implémenté pour attraper les erreurs synchrones et asynchrones, les logger, et renvoyer des réponses d'erreur standardisées au client (avec des codes HTTP appropriés et des messages clairs).

\subsubsection{Points d’Optimisation et Gestion des Performances}
\begin{itemize}[label=\textcolor{HousyGreen}{\textbullet}]
    \item \textbf{Requêtes base de données optimisées} : Indexation des colonnes fréquemment interrogées, utilisation judicieuse des jointures, évitement du problème N+1 avec l'ORM.
    \item \textbf{Caching} : Mise en cache des données fréquemment accédées et peu volatiles (ex: avec Redis).
    \item \textbf{Asynchronisme de Node.js} : Utiliser pleinement la nature non bloquante de Node.js pour gérer efficacement les opérations I/O.
    \item \textbf{Clustering/Load Balancing} : Pour les applications à haute charge, utiliser le module `cluster` de Node.js ou un load balancer (Nginx) pour répartir la charge sur plusieurs processus/serveurs.
    \item \textbf{Compression des réponses HTTP} : Utilisation de `compression` middleware.
    \item \textbf{Monitoring des performances} : Outils comme PM2, New Relic, Datadog pour identifier les goulots d'étranglement.
\end{itemize}
\textit{Un exemple de code montrant un middleware d'authentification ou de gestion d'erreur pourrait être placé ici.}

\subsection{Base de Données et Modélisation}
\label{ssec:bdd_modelisation}
La base de données est le réceptacle de toutes les informations persistantes de Housy. Le choix s'est porté sur PostgreSQL, une base de données relationnelle open-source puissante et fiable.

\subsubsection{Type de BDD, Schéma MER/MLD, Tables Principales et Relations}
Housy utilise une base de données \textbf{PostgreSQL}. La modélisation des données a commencé par un Modèle Entité-Relation (MER) ou Conceptuel de Données (MCD), traduit ensuite en Modèle Logique de Données (MLD) puis en schéma physique SQL.

\begin{figure}[h!]
    \centering
    \includegraphics[width=0.9\linewidth]{./images/placeholder_mer_mld.png}
    \caption{Diagramme Entité-Relation (MER) simplifié de la base de données Housy.}
    \label{fig:mer_housy}
\end{figure}
\textit{Le diagramme MER ou MLD de la base de données Housy, montrant les entités principales et leurs relations, doit être inséré ici.}

\textbf{Tables principales (exemples) :}
\begin{itemize}
    \item `Users` (ou `Utilisateurs`): id (PK), email, password\_hash, nom, prénom, role, created\_at, updated\_at.
    \item `Annonces`: id (PK), titre, description, prix, surface, type\_bien, adresse, ville, code\_postal, user\_id (FK vers Users), created\_at, updated\_at, status (publiée, brouillon, vendue).
    \item `ImagesAnnonce`: id (PK), annonce\_id (FK vers Annonces), url\_image, is\_thumbnail.
    \item `Favoris`: user\_id (FK vers Users), annonce\_id (FK vers Annonces), (PK composite sur user\_id, annonce\_id).
    \item `Messages`: id (PK), sender\_id (FK vers Users), receiver\_id (FK vers Users), annonce\_id (FK vers Annonces, optionnel), content, sent\_at.
\end{itemize}
Les relations sont gérées par des clés étrangères (FK) assurant l'intégrité référentielle. Par exemple, une annonce appartient à un utilisateur (`Annonces.user_id` -> `Users.id`).

\begin{infobox}{ORM Sequelize}
L'ORM (Object-Relational Mapper) Sequelize est utilisé côté backend pour interagir avec la base de données PostgreSQL. Il permet de définir les modèles de données sous forme d'objets JavaScript/TypeScript, de gérer les migrations de schéma, et d'effectuer des opérations CRUD de manière abstraite sans écrire de SQL brut (bien que cela reste possible si besoin).
\begin{codeblock}{Exemple de définition de modèle Sequelize}
// models/annonce.model.ts
import { DataTypes, Model, Sequelize } from 'sequelize';

export class Annonce extends Model {
  public id!: number;
  public titre!: string;
  // ... autres champs
}

export function initAnnonceModel(sequelize: Sequelize): void {
  Annonce.init({
    id: { type: DataTypes.INTEGER, autoIncrement: true, primaryKey: true },
    titre: { type: DataTypes.STRING, allowNull: false },
    // ...
  }, { sequelize, modelName: 'Annonce', tableName: 'annonces' });
}
\end{codeblock}
\end{infobox}

\subsubsection{Requêtes Complexes et Gestion des Migrations}
Pour des besoins spécifiques (recherche avancée, rapports), des requêtes SQL plus complexes peuvent être nécessaires. Sequelize permet d'exécuter du SQL brut ou d'utiliser son API de requêtage avancée (jointures, agrégations, group by).
\textbf{Exemple de requête complexe (conceptuel) :}
"Trouver toutes les annonces de type 'Appartement' à 'Paris', avec un prix entre 1000€ et 2000€, triées par date de création décroissante, incluant les informations de l'utilisateur vendeur et l'image principale."

La \textbf{gestion des migrations} de schéma est cruciale pour l'évolution de la base de données en environnement de développement et de production. Sequelize-CLI fournit des outils pour créer et exécuter des fichiers de migration. Chaque migration contient des fonctions `up` (pour appliquer les changements) et `down` (pour les annuler).
\begin{codeblock}{Exemple de fichier de migration Sequelize}
// migrations/XXXXXXXXXXXXXX-create-annonces.js
'use strict';
module.exports = {
  async up(queryInterface, Sequelize) {
    await queryInterface.createTable('annonces', {
      id: { allowNull: false, autoIncrement: true, primaryKey: true, type: Sequelize.INTEGER },
      titre: { type: Sequelize.STRING, allowNull: false },
      // ... autres colonnes
      userId: {
        type: Sequelize.INTEGER,
        allowNull: false,
        references: { model: 'users', key: 'id' },
        onUpdate: 'CASCADE',
        onDelete: 'CASCADE'
      },
      createdAt: { allowNull: false, type: Sequelize.DATE },
      updatedAt: { allowNull: false, type: Sequelize.DATE }
    });
  },
  async down(queryInterface, Sequelize) {
    await queryInterface.dropTable('annonces');
  }
};
\end{codeblock}

\subsubsection{Sécurité et Optimisation des Données}
\textbf{Sécurité de la Base de Données :}
\begin{itemize}[label=\textcolor{HousyOrange}{\textbullet}]
    \item \textbf{Principe du moindre privilège} : L'utilisateur applicatif de la base de données ne doit avoir que les permissions strictement nécessaires.
    \item \textbf{Prévention des injections SQL} : L'utilisation d'un ORM comme Sequelize aide grandement, car il prépare les requêtes. Si du SQL brut est utilisé, toujours utiliser des requêtes paramétrées.
    \item \textbf{Sauvegardes régulières} : Et tests de restauration.
    \item \textbf{Données sensibles} : Chiffrement des données sensibles au repos (ex: PII) si nécessaire, et toujours en transit (SSL/TLS pour la connexion à la BDD).
\end{itemize}

\textbf{Optimisation des Données :}
\begin{itemize}[label=\textcolor{HousyGreen}{\textbullet}]
    \item \textbf{Indexation} : Créer des index sur les clés étrangères et les colonnes fréquemment utilisées dans les clauses `WHERE`, `JOIN`, `ORDER BY`.
    \item \textbf{Normalisation} : Pour éviter la redondance des données, mais considérer la dénormalisation pour des raisons de performance sur des requêtes de lecture fréquentes.
    \item \textbf{Choix des types de données} : Utiliser les types les plus appropriés et les plus petits possibles pour chaque colonne.
    \item \textbf{Nettoyage et archivage} : Des données anciennes ou inutilisées.
    \item \textbf{Analyse des requêtes} : Utiliser `EXPLAIN ANALYZE` de PostgreSQL pour comprendre et optimiser les plans d'exécution des requêtes lentes.
\end{itemize}
\textit{Un exemple de requête SQL optimisée ou une capture de `EXPLAIN ANALYZE` pourrait être pertinent ici.}

\subsection{Analyse Technologique Croisée}
\label{ssec:analyse_croisee}
Cette section évalue de manière critique la stack technique choisie pour Housy, ses points forts, ses limites, et la compare brièvement à d'autres approches architecturales.

\begin{tipbox}{Points Forts de la Stack Housy}
\begin{itemize}
    \item \textbf{Écosystème JavaScript/TypeScript complet} : Cohérence technologique entre frontend et backend, facilitant le partage de code (ex: types, utilitaires de validation) et la polyvalence des développeurs.
    \item \textbf{Performance de Node.js} : Bien adapté pour les applications I/O-bound et temps réel (si WebSockets utilisés).
    \item \textbf{Popularité et Communauté} : React, Node.js, Express, PostgreSQL bénéficient d'une vaste communauté, de nombreuses bibliothèques, et d'une documentation abondante.
    \item \textbf{Modularité et Scalabilité} : L'architecture choisie permet une bonne séparation des préoccupations, facilitant la maintenance et l'évolution. Le backend peut être scalé horizontalement.
    \item \textbf{Flexibilité de PostgreSQL} : Supporte des données relationnelles et non relationnelles (JSONB), offrant une grande flexibilité.
\end{itemize}
\end{tipbox}

\begin{infobox}{Limites et Points de Vigilance}
\begin{itemize}
    \item \textbf{Nature Monolithique (potentielle) du Backend} : Si le backend Express.js grossit considérablement sans une discipline stricte de modularisation, il peut devenir difficile à maintenir. Une transition vers une architecture microservices pourrait être envisagée à très grande échelle, mais complexifie le développement et le déploiement.
    \item \textbf{Gestion de la Charge CPU Intensive} : Node.js est mono-thread (pour le code utilisateur). Les tâches très CPU-intensives peuvent bloquer la boucle d'événements. Des solutions existent (worker threads, services séparés en Python/Java/Go).
    \item \textbf{Complexité de la Gestion d'État Frontend} : Pour les applications React complexes, la gestion d'état (même avec Redux) peut devenir un défi.
    \item \textbf{Dépendance à l'écosystème NPM} : La gestion des dépendances et les risques de sécurité liés aux packages tiers nécessitent une vigilance constante.
\end{itemize}
\end{infobox}

\subsubsection{Comparaison avec d'autres Architectures Possibles}
\begin{itemize}
    \item \textbf{Backend en Python (Django/Flask) ou Java (Spring Boot)} :
        \begin{itemize}
            \item \textit{Avantages} : Écosystèmes matures, très bonnes performances pour les tâches CPU-bound (surtout Java), outils robustes pour les grandes entreprises. Python est excellent pour l'IA/ML.
            \item \textit{Inconvénients} : Nécessite des compétences différentes du frontend JavaScript. Peut être plus verbeux ou plus lourd (surtout Java Spring).
        \end{itemize}
    \item \textbf{Architecture Microservices} :
        \begin{itemize}
            \item \textit{Avantages} : Scalabilité indépendante des services, choix technologiques par service, résilience améliorée (une panne de service n'impacte pas tout).
            \item \textit{Inconvénients} : Complexité de gestion (déploiement, communication inter-services, monitoring, transactions distribuées), surcoût initial. Adapté pour les très grandes applications.
        \end{itemize}
    \item \textbf{Backend as a Service (BaaS) / Serverless (AWS Lambda, Firebase)} :
        \begin{itemize}
            \item \textit{Avantages} : Réduction de la gestion d'infrastructure, scalabilité automatique, potentiellement moins cher pour des trafics variables.
            \item \textit{Inconvénients} : Vendor lock-in, limitations sur l'exécution (temps, mémoire), complexité pour la logique métier très élaborée, "cold starts".
        \end{itemize}
    \item \textbf{Frontend avec un autre framework (Vue.js, Angular)} :
        \begin{itemize}
            \item \textit{Vue.js} : Souvent considéré comme plus facile à prendre en main que React, très performant.
            \item \textit{Angular} : Framework complet et opinioné, très structuré, souvent privilégié pour les grandes applications d'entreprise. TypeScript est natif.
            \item Le choix de React pour Housy était pertinent compte tenu de sa popularité et de son écosystème.
        \end{itemize}
\end{itemize}
Le choix de la stack actuelle pour Housy représente un bon compromis entre performance, productivité de développement, et scalabilité pour les besoins du projet tels que définis.

\subsection{Illustrations et Captures d’Application}
\label{ssec:illustrations_captures}
Cette section est dédiée à la présentation visuelle de l'application Housy. Des captures d'écran permettent d'illustrer concrètement les fonctionnalités et l'interface utilisateur.

\begin{figure}[h!]
    \centering
    \includegraphics[width=0.7\linewidth]{./images/placeholder_homepage.png}
    \caption{Capture d'écran de la page d'accueil de Housy.}
    \label{fig:homepage}
\end{figure}
\textit{Insérer ici une capture d'écran de la page d'accueil de Housy.}

\begin{figure}[h!]
    \centering
    \includegraphics[width=0.7\linewidth]{./images/placeholder_annonce_detail.png}
    \caption{Capture d'écran de la page de détail d'une annonce sur Housy.}
    \label{fig:annonce_detail}
\end{figure}
\textit{Insérer ici une capture d'écran montrant la page de détail d'une annonce.}

\begin{figure}[h!]
    \centering
    \includegraphics[width=0.7\linewidth]{./images/placeholder_dashboard_admin.png}
    \caption{Capture d'écran du tableau de bord administrateur (si applicable).}
    \label{fig:dashboard_admin}
\end{figure}
\textit{Insérer ici une capture d'écran d'une interface d'administration ou d'un dashboard pertinent.}

\begin{figure}[h!]
    \centering
    \includegraphics[width=0.7\linewidth]{./images/placeholder_formulaire_creation.png}
    \caption{Capture d'écran d'un formulaire de création (ex: création d'annonce).}
    \label{fig:formulaire_creation}
\end{figure}
\textit{Insérer ici une capture d'écran d'un formulaire clé de l'application.}

Chaque capture d'écran doit être accompagnée d'une légende claire expliquant ce qui est montré et, si pertinent, mettant en évidence des aspects techniques ou fonctionnels spécifiques.

\subsection{Résumé et Recommandations Techniques}
\label{ssec:resume_recommandations}
Le développement de Housy a permis de mettre en œuvre une application web fullstack moderne et fonctionnelle. L'architecture choisie, basée sur React pour le frontend et Node.js/Express pour le backend avec une base de données PostgreSQL, s'est avérée robuste et flexible.

\begin{tipbox}{Bonnes Pratiques Appliquées et Retenues}
\begin{itemize}
    \item \textbf{Séparation des préoccupations} : Clairement définie entre frontend, backend, et base de données, ainsi qu'au sein de chaque composant (ex: couches du backend, composants React).
    \item \textbf{Utilisation de TypeScript} : A grandement amélioré la qualité du code, la détection précoce des erreurs et la maintenabilité.
    \item \textbf{Tests automatisés} : Essentiels pour garantir la non-régression et la fiabilité (tests unitaires, d'intégration).
    \item \textbf{Gestion de version avec Git} : Suivi rigoureux des branches, conventions de commit.
    \item \textbf{CI/CD (Intégration Continue / Déploiement Continu)} : Automatisation des builds, tests et déploiements pour accélérer les livraisons et réduire les erreurs.
    \item \textbf{Principes SOLID} : Application (autant que possible) des principes SOLID dans la conception logicielle, en particulier dans le backend.
\end{itemize}
\end{tipbox}

\begin{infobox}{Points de Vigilance et Recommandations Futures}
\begin{itemize}
    \item \textbf{Scalabilité à long terme} : Surveiller les performances et envisager des optimisations plus poussées (caching avancé, réplication de base de données, passage à des microservices si la complexité l'exige) si l'application connaît une forte croissance.
    \item \textbf{Sécurité continue} : La sécurité n'est pas un état mais un processus. Mettre en place des audits réguliers, veille sur les vulnérabilités des dépendances (NPM Audit, Snyk).
    \item \textbf{Monitoring et Alerting} : Déployer des outils de monitoring complets (performances applicatives, erreurs, logs) avec des alertes pour réagir proactivement aux problèmes.
    \item \textbf{Documentation technique} : Maintenir à jour la documentation de l'API (Swagger/OpenAPI), les diagrammes d'architecture, et les décisions de conception.
    \item \textbf{Tests d'utilisabilité et retours utilisateurs} : Continuer à recueillir les retours pour améliorer l'UX et prioriser les futures fonctionnalités.
    \item \textbf{Optimisation des coûts d'infrastructure} : Si hébergé sur le cloud, surveiller et optimiser les coûts (choix des instances, auto-scaling, services managés).
    \item \textbf{Formation continue de l'équipe} : Les technologies évoluent vite, encourager la veille technologique et la formation.
\end{itemize}
\end{infobox}

En conclusion, le projet Housy repose sur des fondations techniques solides. Les choix effectués offrent un bon équilibre pour les besoins actuels et une base saine pour les évolutions futures. L'accent mis sur la qualité du code, la sécurité et les bonnes pratiques de développement est un gage de pérennité pour l'application.

% Fin du chapitre
\end{document}
